% ===== 抑制 Quarto 默认 \maketitle（自定义封面代行）=====
\let\maketitle\relax
\AtBeginDocument{\let\maketitle\relax}
% ===== 粗体数学 =====
% 注意：不要加载 bm 宏包！bm 会重定义 \boldsymbol，且与 unicode-math 的
% OpenType 数学字体不兼容，编译时报 "Improper alphabetic constant"。
% 统一使用 amsmath 的 \boldsymbol{X}（KaTeX/MathJax/xelatex 三方兼容）。
% ===== 中文段落：段首空两字、统一设置 =====
\setlength{\parindent}{2em}    % 段首缩进 2 个汉字宽
\setlength{\parskip}{0pt}      % 段间不留额外垂直空白（中文书传统）
\usepackage{indentfirst}       % 让 \chapter 后的第一段也跟着缩进
\usepackage[a4paper, inner=2.5cm, outer=2.0cm, top=2.5cm, bottom=2.5cm]{geometry}
\usepackage{setspace}
\setstretch{1.35}
% ===== 书籍样式页眉页脚 =====
\usepackage{fancyhdr}
\usepackage{xcolor}
\usepackage{titlesec}
\pagestyle{fancy}
\usepackage{amsthm}
% 公理/定义环境已改用 Quarto 原生定理 div（::: {#prp-*} / ::: {#def-*}），
% 其 LaTeX 环境由 Quarto 自动定义，此处不再 \newtheorem，避免重复定义冲突。
\fancyhf{}
\fancyhead[LE]{\thepage\quad\textit{狭义相对论}}
\fancyhead[RO]{\textit{狭义相对论}\quad\thepage}
\fancyfoot[C]{\thepage}
\renewcommand{\headrulewidth}{0.4pt}
\renewcommand{\footrulewidth}{0pt}
\fancypagestyle{plain}{
    \fancyhf{}
    \fancyfoot[C]{\thepage}
    \renewcommand{\headrulewidth}{0pt}
    \renewcommand{\footrulewidth}{0pt}
}
% ===== 章节标题美化 =====
% 强制全局使用阿拉伯数字编号（第1章）
\CTEXsetup[
    name={第,章},
    number=\arabic{chapter}
]{chapter}
% 章号自动改为"第N章"（ctexbook 内置），并按 display 样式居中。
\titleformat{\chapter}
    {\Huge\bfseries\color{blue!50!black}\centering}
    {第\thechapter 章}{1em}{}

\titleformat{\section}
    {\Large\bfseries\color{blue!40!black}}
    {\thesection}{1em}{}

\titleformat{\subsection}
    {\large\bfseries\color{blue!30!black}}
    {\thesubsection}{0.8em}{}

\titleformat{\subsubsection}
    {\normalsize\bfseries\color{blue!20!black}}
    {\thesubsubsection}{0.6em}{}

% ===== 5、6 级标题（段落头 + runin）：必须补齐，
%        否则一旦文档里出现 ##### ###### 就会报
%        "Package titlesec Error: No format for this command." =====
\titleformat{\paragraph}[runin]
    {\normalsize\bfseries}{}{0pt}{}
\titleformat{\subparagraph}[runin]
    {\normalsize\itshape}{}{0pt}{}
